% English translation of chapter1.tex, lines 9--31.
% Source: Methods of Algebra, Volume 2, commit
% 9a5803ff77dd3257484cb177f851a73770a59dd3 (CC BY 4.0).
% stable-unit-id: o014.aljabr2.chapter1.overview

% segment-id: o014.li2.chapter1.overview.h001
\chapter{Supplements to Category Theory}\label{sec:categories}
\hypertarget{o014.aljabr2.chapter1.overview}{ }

% segment-id: o014.li2.chapter1.overview.p001
As its title suggests, this chapter introduces several tools from category
theory that will be needed in later chapters, both as a review and as an
opportunity to see familiar ideas afresh. Some of this material appears in
\cite{Li1} or similar textbooks; it is nevertheless restated here in order to
review the concepts and standardize the notation. For other basic terminology,
the reader may consult the introduction to this book.

% segment-id: o014.li2.chapter1.overview.l001
\begin{itemize}
	% segment-id: o014.li2.chapter1.overview.i001
	\item Sections~\sourcecrossref{sec:subquot}{1.1}--%
	\sourcecrossref{sec:images}{1.2} begin with a brief introduction to
	subobjects and quotient objects, including the general properties of
	monomorphisms and epimorphisms. They then introduce the image $\Image(f)$
	and coimage $\Coim(f)$ of a morphism $f$, as well as morphisms for which
	``image equals coimage,'' called strict morphisms. These are basic notions
	that make sense in arbitrary categories.

	% segment-id: o014.li2.chapter1.overview.i002
	\item Sections~\sourcecrossref{sec:Ker-Coker}{1.3}--%
	\sourcecrossref{sec:linear-cat}{1.4} form the linear part of the chapter.
	They introduce a special kind of category in which every morphism set has
	the structure of an abelian group, called an $\cate{Ab}$-category; a functor
	that preserves addition of morphisms is called additive. In such categories
	one can define kernels and cokernels of morphisms, generalizing those of
	module homomorphisms. An $\cate{Ab}$-category with a zero object and finite
	products is called an additive category. In an additive category one can
	form finite direct sums of objects, and morphisms between finite direct sums
	can be manipulated in much the same way as matrices.

	More generally, the additive structure can be enhanced with scalar
	multiplication by a commutative ring $\Bbbk$, leading to the notions of a
	$\Bbbk\dcate{Mod}$-category and a $\Bbbk$-linear category. Categories with
	these linear structures are a principal concern of this book.

	% segment-id: o014.li2.chapter1.overview.i003
	\item Sections~\sourcecrossref{sec:limit-functor}{1.5}--%
	\sourcecrossref{sec:filtered-indlim}{1.6} cover the foundations of limits.
	In particular, Definition~\sourcecrossref{def:create-limit}{1.5.2}
	introduces the terminology for a functor $F:\mathcal C\to\mathcal D$
	that preserves limits (from $\mathcal C$ to $\mathcal D$) or creates them
	(from $\mathcal D$ to $\mathcal C$). Cofinal functors and filtered
	inductive limits are also discussed; these are useful and important
	notions.

	% segment-id: o014.li2.chapter1.overview.i004
	\item Left and right Kan extensions, discussed in Sections~\sourcecrossref{sec:Kan-extension}{1.7}--%
	\sourcecrossref{sec:Kan-as-limit}{1.8}, are fundamental operations in
	category theory. Under suitable conditions they can be constructed using
	$\varinjlim$ or $\varprojlim$. The argument is somewhat involved, but the
	underlying idea is simple: express the problem of extending a functor as a
	``best approximation'' formulated in terms of limits. We also record the
	non-obvious Theorem~\sourcecrossref{prop:Quillen-adjunction-gen}{1.7.7},
	which explains the connection between absolute left or right Kan extensions
	and adjoint pairs.

	% segment-id: o014.li2.chapter1.overview.i005
	\item Gabriel--Zisman localization, introduced in
	Section~\sourcecrossref{sec:cat-localization}{1.9} and henceforth called
	simply \emph{localization}, is the construction that formally adjoins
	inverses to a family $S$ of morphisms in a category. It is analogous to
	localization of rings, but the concrete details are considerably more
	complicated. To retain better control over the resulting category, in
	practice $S$ is often assumed to be a \emph{multiplicative system} in the
	sense of Definition~\sourcecrossref{def:multiplicative-family}{1.9.5}.
	The principal use of localization in this book is to construct derived
	categories; another is to define Serre quotients of abelian categories in
	Section~\sourcecrossref{sec:Serre-subcat}{2.9}.
	Section~\sourcecrossref{sec:functor-localization}{1.10} explains how to
	form Kan extensions along a localization of categories. This may be viewed
	as the localization of a functor and is essential for understanding derived
	functors.

	% segment-id: o014.li2.chapter1.overview.i006
	\item The adjoint functor theorem discussed in
	Section~\sourcecrossref{sec:AFT}{1.11} characterizes, under suitable
	conditions, when a functor has a left or right adjoint. Its principal
	application in this book is to the study of Grothendieck categories in
	Section~\sourcecrossref{sec:Grothendieck-cat}{2.10}. The theorem has uses
	far beyond this application, so the section is valuable in its own right.
	The notions of generator and cogenerator discussed there
	(Definition~\sourcecrossref{def:generators}{1.11.8}) are also essential.
\end{itemize}

% segment-id: o014.li2.chapter1.overview.p002
This chapter draws extensively on \cite{KS06, Rie16}.

% segment-id: o014.li2.chapter1.overview.n001
\begin{petunjukbacaan}
The basic theory of abelian categories and complexes requires only
Sections~\sourcecrossref{sec:subquot}{1.1}--%
\sourcecrossref{sec:linear-cat}{1.4} of this chapter.
Sections~\sourcecrossref{sec:limit-functor}{1.5}--%
\sourcecrossref{sec:filtered-indlim}{1.6} provide supporting material and are
also cited repeatedly in other chapters. The theory of derived categories uses
Sections~\sourcecrossref{sec:Kan-extension}{1.7}--%
\sourcecrossref{sec:functor-localization}{1.10}. This material has applications
beyond that theory, but readers may judge their own needs and postpone it until
just before studying derived categories. The final section,
Section~\sourcecrossref{sec:AFT}{1.11}, is used only in
Section~\sourcecrossref{sec:Grothendieck-cat}{2.10}; nevertheless, its material
on generators and cogenerators is essential and may be consulted as needed.
\end{petunjukbacaan}
